\documentclass[../../../main.tex]{subfiles}
\begin{document}

\subsection{Meta-Regression}

Usually, regression models are based on data comprising individual persons or specimens, for which both the independent $x$ and dependent $y$ values are measured. In meta-regression, this concept is applied to entire studies. The variable $x$ represents characteristics of studies, for example, the year in which it was conducted. Based on this information, a meta-regression model tries to predict $y$, the study's effect size \autocite{Harrer2021}.

In a conventional regression, we want to estimate the value $y_i$ of a person or specimen $i$ using a predictor (or covariate) $x_i$ with a regression coefficient $\beta$ \autocite{Harrer2021}.

\begin{equation}
y_i = \beta_0 + \beta_1 x_i.
\end{equation}

In meta-regression, the variable $y$ we want to predict is the observed effect size $\hat{\theta}_k$ of study $k$ \autocite{Harrer2021}.

\begin{equation}
\hat{\theta}_k = \theta + \beta x_k + \epsilon_k + \zeta_k.
\label{Equation:MetaRegression}
\end{equation}

This formula contains two extra terms, $\epsilon_k$ and $\zeta_k$. The same terms found in the random-effects model from \autoref{Equation:RandomEffects}. The first one, $\epsilon_k$, is the sampling error through which the effect size of a study deviates from its true effect. The second error, $\zeta_k$, denotes that even the true effect size of the study is only sampled from an overarching distribution of effect sizes. Since the meta-regression formula includes a fixed effect (the $\beta$ coefficient) as well as a random effect ($\zeta_k$), the model used in meta-regression is often called a mixed-effects model \autocite{Harrer2021}.

\subsubsection{Meta-Regression with a Categorical Predictor}

Subgroup analysis (from \autoref{Subsection:SubgroupAnalysis}) is nothing else than a meta-regression with a categorical predictor, where such categorical variables can be included through dummy-coding, such as the following \autocite{Harrer2021}.

\begin{equation}
D_g =
\begin{cases}
0: &\text{Subgroup A} \\
1: &\text{Subgroup B} \\
\end{cases}
\end{equation}

To specify a subgroup analysis in the form of a meta-regression, we simply have to replace the covariate $x_k$ with $D_g$ \autocite{Harrer2021}.

\begin{equation}
\hat{\theta}_k = \theta + \beta D_g + \epsilon_k + \zeta_k.
\end{equation}

\begin{figure}[!htbp]
\centering
\includegraphics[width=0.8\linewidth]{appendix/methods/analysis/assets/ma/subgroup_regression.png}
\caption[Meta-regression]{Meta-regression with a categorical predictor (subgroup analysis).}
\end{figure}

\subsubsection{Meta-Regression with a Continuous Predictor}

When people speak of a meta-regression, they usually think of models in which a continuous variable is used as the predictor. The aim of the meta-regression model is to find values of $\theta$ and $\beta$ which minimize the difference between the predicted effect size and the true effect size of studies. Taking into account both the sampling error $\epsilon_k$ and between-study heterogeneity $\zeta_k$, meta-regression tried to find a model that generalizes well, not only to the observed effect sizes but to the universe of all possible studies of interest \autocite{Harrer2021}.

An important detail about meta-regression models is that they can be seen as an extension of the random-effects model. The random-effects model is nothing but a meta-regression without a slope term. Since it contains no slope, the random-effects model simply predicts the same value for all studies. That is, the estimate of the pooled effect size $\mu$, which is equivalent to the intercept \autocite{Harrer2021}.

\begin{figure}[!htbp]
\centering
\begin{subfigure}{0.37\textwidth}
\includegraphics[width=\linewidth]{appendix/methods/analysis/assets/ma/regression_left.png}
\end{subfigure}
\hspace{3em}
\begin{subfigure}{0.37\textwidth}
\includegraphics[width=\linewidth]{appendix/methods/analysis/assets/ma/regression_right.png}
\end{subfigure}
\caption[Random-effects model and meta-regression]{Example of a random-effects model versus a meta-regression.}
\end{figure}

In the first step, the calculation of a meta-regression therefore closely resembles the one of a random-effects meta-analysis, in that the between-study heterogeneity $\tau^2$ is estimated. In the next step, the fixed weights $\theta$ and $\beta$ are estimated. In meta-regression, weighted least squares (WLS) is used to find the regression line that fits the data best. WLS makes sure that studies with a smaller standard error are given a higher weight \autocite{Harrer2021}.

If the meta-regression model fits the data well, the true effect sizes should deviate less from the regression line compared to the pooled effect $\hat{\mu}$. If this is the case, the predictor $x$ explains some of the heterogeneity variance in our meta-analysis. The fit of the meta-regression model can therefore be assessed by checking how much of the heterogeneity variance it explains \autocite{Harrer2021}.

The predictors included in the mixed-effects model should minimize the amount of residual heterogeneity variance denoted by $\hat{\tau}_\text{MEM}^2$. The $R_\ast^2$ index is used to quantify the percentage of variation explained by the meta-regression model. The asterisk is added to differentiate from the $R^2$ index used in conventional regressions, since meta-regression deals with true effect sizes instead of observed data points \autocite{Harrer2021}.

$R_\ast^2$ uses the amount of residual heterogeneity variance that even the meta-regression slope cannot explain. It is put in relation to the total heterogeneity that we initially found in our random-effects model $\hat{\tau}_\text{REM}^2$. Subtracting this fraction from 1 leaves us with the percentage of between-study heterogeneity explained by the predictor, and is expressed in \autoref{Equation:RegressionIndex}. We can also says that it is how much the mixed-effects model has reduced the heterogeneity variance compared to the initial random-effects pooling model, in percent \autocite{Harrer2021}.

\begin{equation}
R_\ast^2
= 1 - \frac{\hat{\tau}_\text{MEM}^2}{\hat{\tau}_\text{REM}^2}
= \frac{\hat{\tau}_\text{REM}^2-\hat{\tau}_\text{MEM}^2}{\hat{\tau}_\text{REM}^2}.
\label{Equation:RegressionIndex}
\end{equation}

Usually, we are not only in the amount of heterogeneity explained by the regression model but also if the regression weight of our predictor $x$ is significant. Both in conventional and meta-regression, the significance of a regression weight is commonly assessed through a Wald test. This involves calculating the test statistic $z$, by dividing the estimate of $\beta$ through its standard error \autocite{Harrer2021}.

\begin{equation}
z = \frac{\beta}{SE_{\hat{\beta}}}.
\end{equation}

Under the null hypothesis that $\beta=0$, this $z$-statistic follows a standard normal distribution. This allows us to calculate a corresponding $p$-value. Like in normal meta-analysis models, we can also use the Knapp-Hartung adjustment, which results in a test statistic based on the $t$-distribution \autocite{Harrer2021}.
\end{document}
